\chapter*{Annexes}
\addcontentsline{toc}{chapter}{Annexes}

\section*{Annexe A : Organisation des diagrammes UML}

Les diagrammes UML du projet sont regroupés dans le dossier \texttt{diagrammes/}. Chaque fichier de sprint contient plusieurs pages Draw.io afin de garder dans un même livrable les vues de conception d'une même itération.

\begin{longtable}{p{4cm} p{9cm}}
\toprule
\textbf{Fichier} & \textbf{Contenu} \\
\midrule
\texttt{classe\_global.drawio} & Diagramme de classes global et diagramme de cas d'utilisation global. \\
\texttt{sprint1.drawio} & Classes, cas d'utilisation, séquence d'authentification JWT et séquence de résolution tenant. \\
\texttt{sprint2.drawio} & Classes, cas d'utilisation, séquence de lecture dynamique et séquence d'importation intelligente. \\
\texttt{sprint3.drawio} & Classes, cas d'utilisation, séquence de sauvegarde de template et séquence de prévisualisation avec fusion. \\
\texttt{sprint4.drawio} & Classes, cas d'utilisation, séquence de génération en masse et séquence d'audit soft-delete. \\
\bottomrule
\end{longtable}

\section*{Annexe B : Captures d'écran à intégrer}

Les captures suivantes peuvent être ajoutées dans le dossier \texttt{images/} au moment de la préparation de la version finale imprimable :

\begin{itemize}
    \item interface de connexion et redirection selon le rôle ;
    \item écran de sélection de l'année universitaire ;
    \item tableau de bord administrateur ;
    \item configuration d'un module métier et d'une TableView ;
    \item configuration d'une famille documentaire ;
    \item studio d'édition d'un template ;
    \item prévisualisation d'un document fusionné ;
    \item écran de génération documentaire ;
    \item liste des archives avec filtres ;
    \item détail d'une archive et affichage du document généré.
\end{itemize}

\section*{Annexe C : Exemples de composants techniques}

Les éléments suivants constituent des points d'entrée utiles pour relier le rapport au code source :

\begin{itemize}
    \item \texttt{backend/Middlewares/TenantResolutionMiddleware.cs} pour la résolution du tenant ;
    \item \texttt{backend/Common/Tenant/TenantResolver.cs} pour la lecture de l'organisation ;
    \item \texttt{backend/Repositories/EditorRepository.cs} pour l'accès SQL dynamique ;
    \item \texttt{frontend/src/app/features/editor/services/document-render.service.ts} pour le rendu documentaire ;
    \item \texttt{frontend/src/app/features/editor/services/table-view.service.ts} pour la gestion des vues dynamiques ;
    \item \texttt{frontend/src/app/features/editor/pages/document-archive/} pour la consultation des archives.
\end{itemize}
